%=====================================================================
%  译校说明（Translation Notes）
%  本文件为译者所加，非原文内容，置于参考文献之后。
%=====================================================================
\section*{译校说明}
\addcontentsline{toc}{section}{译校说明}

本译文由英文学术论文 \emph{Supersymmetric field theories and geometric Langlands: The other side of the coin}
（A.~Balasubramanian, J.~Teschner, arXiv:1702.06499v2，提交于 2017 年 3 月 10 日，String-Math 2016 会议文集）
精译而成。为便于读者核对与引用，此处说明翻译所依据的底本、处理原则与已知问题。

\subsection*{一、底本与工作方式}

\begin{enumerate}[label=（\arabic*）,leftmargin=*]
  \item \textbf{底本}：以 arXiv 公开的 \LaTeX{} 源码（\texttt{AGT-GL.tex}, 2320 行，及 \texttt{AGT-GL.bbl}）为主要依据，
        并与随任务提供的 PDF（35 页，经逐页核对与 arXiv 版一致）对照。以源码而非 OCR 为准，
        可保证公式、指标、上下标、希腊字母在转换中绝对无损。
  \item \textbf{公式}：所有行内公式与显示公式均逐字保留源码原貌，包括 \verb|\label|、\verb|\left/\right|、\verb|\,|、
        矩阵环境与 \verb|\underset| 等排版细节；未重排、未化简、未“顺手”改符号。原文中被作者注释掉的公式
        （如附录中 Hitchin 方程的自对偶形式）保持注释状态，仅在批注中说明其内容。
  \item \textbf{正文}：按中文学术论文的语序与行文习惯重写，不作逐词硬译；术语以文末《术语译名对照表》为准，
        人名一律保留拉丁字母。
  \item \textbf{批注}：正文源码中以 \verb|% [Tutor Note: ...]| 形式插入译者批注，共 44 条，分“易错提醒”与
        “跳步点拨”两类。批注只出现在源码注释中，不影响排版；在 .tex 源码中可直接看到。
        各节开头另有“本节重点（Key Takeaways）”，为译者所加，非原文内容。
  \item \textbf{编译}：本文档使用 XeLaTeX 编译（中文支持依赖 \texttt{ctex} 宏包）。\texttt{pdfLaTeX} 不能直接编译含中文的版本。
\end{enumerate}

\subsection*{二、图表清单}

\begin{enumerate}[label=（\arabic*）,leftmargin=*]
  \item 经核查，原文 PDF 共 35 页，\textbf{不含任何 Figure/Diagram 图像对象}
        （PDF 内部图像对象计数为 0），故无图可裁、无图可移。
  \item 原文含两张表格，均已原样保留于正文并作了表头翻译：
        第 3 节的 M5/M5$'$ 膜支撑方向表，第 5 节的 M5/M5$'$/M2 膜支撑方向表。
        原文未给表题，故未添加 \verb|\caption|。
\end{enumerate}

\subsection*{三、原文可疑之处与已知笔误}

以下各条均为译者在逐式核对中发现的原文问题。译文一律\textbf{保留原样}，另以批注说明；此处集中列出，供教授复核。

\begin{enumerate}[label=\textbf{[\arabic*]},leftmargin=*]
  \item \textbf{记号不统一（多节）}：原文 $N=4$ 与 $\mathcal{N}=4$、$\mathcal{N}=(4,4)$ 混用；
        ${}^{\rm L}G$ 与 ${}^{L}G$ 混用；术语表上标有 \verb|\glos|、$^{g)}$、$^{(g)}$ 三种写法。
  \item \textbf{§2.5 印刷错误}：geometric Lang\textbf{s}lands（应为 Langlands）；
        “branes \textbf{the} dual sigma model” 缺介词 of。
  \item \textbf{§2.5 逻辑疑点}：Cordova--Jafferis 一段称边界条件位于“$C$ 的两端”，但 $C$ 为紧 Riemann 面、无端点；
        按上下文（$I_0\times C$ 上的复 Chern--Simons）疑应为区间 $I_0$ 的两端。
  \item \textbf{§3 式 (3.3)}：正文只定义 $\mathfrak{t}_+=\mathfrak{t}_1+i\mathfrak{t}_2$，式中却用 $\mathfrak{t}_-$；
        第三条的 $\mathcal{A}_1$ 与前文 $A_y$ 记法不统一。
  \item \textbf{§3}：$\mathcal{A}_{\rm cc}^\hbar=\mathrm{Hom}_{\CM_G(C)}(\cdots)$ 中 $\CM_G(C)$ 全文未定义，
        按上下文应为 $\CM_H(G)$；$T^*\mathcal{T}(C)|_{C}$ 的下标应为 Teichmüller 空间中的点 $[C]$。
  \item \textbf{§3}：“very stable 丛 = 不承认幂零 Higgs 场”漏掉“非零”二字（零 Higgs 场总是存在）。
  \item \textbf{§4 指标错位（系统性）}：$\langle\de_a^{\sst(2)},\phi_x\rangle$ 应为 $\de_a^{\sst(1)}$；
        $\vf_a\in\CH_x^{\sst(1)}$ 应为 $\CH_a^{\sst(1)}$；$\mathsf{L}_{\ep,a}^{\sst(1)\si}$ 多一个 $\ep$ 下标；
        配分函数自变量 $(a;x;\tau)$、$(a;x,\tau)$、$(x,a,\tau)$ 三种写法并存。
  \item \textbf{§4 宏误用}：原文 $S^1_{\CS}$ 处的 \verb|\CS| 在导言区定义为花体 $\mathcal{S}$（class $\mathcal{S}$ 的 S），
        但此处语义为 Chern--Simons；译文按“复 CS（Chern--Simons）理论”处理。
  \item \textbf{§6}：式 (6.3) 写 $\mathcal{Y}^{\rm 4d}$，紧接一句却称 $\mathcal{W}^{\rm 4d}$；属笔误。
  \item \textbf{附录 括号不配对}：$\BC[\mathrm{Hom}(\pi_1(C),G]^G$ 内层圆括号未闭合；
        且不变量环 $\BC[\mathrm{Hom}]^G$ 与 GIT 商 $\mathrm{Hom}/\!/G$ 混写。
  \item \textbf{附录 $\mathrm{End}$ 应为 $\mathrm{End}_0$}：$G=SL(2)$ 时 Higgs 场取值应为无迹自同构丛，
        否则与 $\CB\simeq H^0(C,K^2)$ 及 $GL(2)$ 情形不自洽。
  \item \textbf{附录 谱曲线方程与 $GL(2)$ 讨论错配}：$v^2=\frac12\mathrm{tr}(\vf^2)$ 只在 $\mathrm{tr}\,\vf=0$ 时成立；
        紧随其后的 $G=GL(2)$ 讨论中正确方程应为 $\det(v-\vf)=0$，即 $v^2-(\mathrm{tr}\vf)v+\det\vf=0$。
  \item \textbf{附录 退化未注明}：“由 $(\CF,\nabla_\ep')$ 的和乐得到基本群表示”只在 $\ep\neq0$ 时成立；
        $\ep=0$ 时 $\nabla_0'=\vf$ 退化为 Higgs 场，“和乐”无意义。
  \item \textbf{附录 下标混用}：$L_i=\mathrm{tr}(\rho(\ga_i))$ 用 $i$，随后写作 $L_r=2\cosh(a_r/2)$，$i$ 与 $r$ 混用。
\end{enumerate}

\subsection*{四、阅读建议}

\begin{enumerate}[label=（\arabic*）,leftmargin=*]
  \item 建议先读 §2.1（Beilinson--Drinfeld 进路）、§2.2（Kapustin--Witten 进路）、§2.3（Nekrasov--Witten/AGT），
        再读 §2.4 指出的共同根源（六维 $(2,0)$ 理论在 $\CM^6=\CM^4\times C$ 上的两种紧化）。
        本文的核心主张在 §6 与 §7：带面算符的 AGT 对应之 Nekrasov--Shatashvili 极限给出几何 Langlands 对应。
  \item 附录 “Hitchin's moduli spaces” 是全文术语的汇总，建议与术语译名对照表配合使用。
  \item 源码中 44 条 \verb|% [Tutor Note: ...]| 批注按“易错/跳步”标注，可直接用于组会报告或备课。
\end{enumerate}
